\documentclass[aspectratio=169,12pt]{beamer}
% Maths Lycée — thème éditorial inspiré du portfolio de Yassine.
% Modifier ici les couleurs et la typographie de tous les diaporamas.
\usepackage[T1]{fontenc}
\usepackage[utf8]{inputenc}
\usepackage[french]{babel}
\usepackage{lmodern,amsmath,amssymb,mathtools,tikz,pgfplots,listings,adjustbox}
\usetikzlibrary{arrows.meta,positioning,calc}
\pgfplotsset{compat=1.18,width=9cm,height=4.2cm}
\definecolor{Ink}{HTML}{202124}
\definecolor{Blue}{HTML}{356B9E}
\definecolor{Mist}{HTML}{F6F6F5}
\definecolor{Muted}{HTML}{737579}
\definecolor{Line}{HTML}{E7E7E5}
\setbeamercolor{normal text}{fg=Ink,bg=white}
\setbeamercolor{structure}{fg=Blue}
\setbeamercolor{frametitle}{fg=Ink}
\setbeamercolor{title}{fg=Ink}
\setbeamercolor{block title}{fg=Blue,bg=Mist}
\setbeamercolor{block body}{fg=Ink,bg=Mist}
\setbeamerfont{title}{size=\LARGE,series=\bfseries}
\setbeamerfont{frametitle}{size=\Large,series=\bfseries}
\setbeamersize{text margin left=1cm,text margin right=1cm}
\setbeamertemplate{navigation symbols}{}
\setbeamertemplate{itemize items}[circle]
\setbeamertemplate{enumerate items}[default]
\setlength{\parskip}{7pt}
\newcommand{\slidekind}{COMPRENDRE}
\newcommand{\chapternumber}{01}
\setbeamertemplate{frametitle}{\vspace{0.3cm}{\fontsize{7}{9}\selectfont\color{Blue}\MakeUppercase{\slidekind}}\par\vspace{0.1cm}\insertframetitle\par\vspace{0.08cm}\color{Line}\rule{\textwidth}{0.4pt}\vspace{-0.02cm}}
\setbeamertemplate{footline}{\hspace{1cm}\begin{minipage}{\dimexpr\paperwidth-2cm}\color{Line}\rule{\textwidth}{0.4pt}\par\vspace{0.06cm}{\fontsize{6.5}{8}\selectfont\color{Muted}MATHS LYCÉE\quad / \quad PREMIÈRE SPÉCIALITÉ\hfill\insertframenumber\ /\ \inserttotalframenumber}\end{minipage}\vspace{0.22cm}}
\setbeamertemplate{title page}{
\vspace{0.3cm}
{\fontsize{8}{10}\selectfont\color{Blue}MATHS LYCÉE\quad /\quad 2026--2027}\par
\vspace{0.6cm}
\begin{columns}[c,onlytextwidth]
\begin{column}{.72\textwidth}
{\fontsize{9}{12}\selectfont\color{Muted}CHAPITRE \chapternumber}\par\vspace{0.22cm}
{\usebeamerfont{title}\color{Ink}\inserttitle\par}
\vspace{0.4cm}{\color{Blue}\rule{1.5cm}{2pt}}\par
{\small\color{Muted}Comprendre les notions.\par Justifier chaque étape.}
\end{column}
\begin{column}{.23\textwidth}
\begin{tikzpicture}[scale=.65]
\fill[Mist,rounded corners=6pt](-1.9,-2.1)rectangle(2,2.5);
\draw[Line,step=.5](-1.8,-2)grid(1.9,2.4);
\draw[Muted,->](-1.7,-1.3)--(1.8,-1.3);
\draw[Muted,->](0,-1.6)--(0,2.2);
\draw[Blue,very thick,domain=-1.5:1.5,samples=40]plot(\x,{1.3*\x*\x-1.3});
\fill[Blue](1,0)circle(2.5pt);
\end{tikzpicture}
\end{column}
\end{columns}
\vfill}
\lstdefinestyle{mathspython}{language=Python,basicstyle=\ttfamily\small,keywordstyle=\color{Blue}\bfseries,commentstyle=\color{Muted},stringstyle=\color{Blue},backgroundcolor=\color{Mist},frame=none,xleftmargin=8pt,xrightmargin=8pt,aboveskip=8pt,belowskip=8pt,showstringspaces=false,columns=fullflexible,keepspaces=true,breaklines=true}
\lstnewenvironment{codepython}{\lstset{style=mathspython}}{}
\newcommand{\R}{\mathbb{R}}
\newcommand{\N}{\mathbb{N}}
\newcommand{\E}{\mathbb{E}}
\newcommand{\V}{\mathrm{V}}
\hypersetup{colorlinks=true,urlcolor=Blue}

\title{Produit scalaire}
\subtitle{Première spécialité}
\renewcommand{\chapternumber}{09}\author{Maths Lycée}
\date{Année scolaire 2026--2027}
\begin{document}\begin{frame}\titlepage\end{frame}
\begin{frame}{Objectifs du chapitre}\begin{itemize}
\item Calculer un produit scalaire avec la formule adaptée
\item Démontrer une orthogonalité
\item Déterminer une longueur ou un angle
\item Utiliser Al-Kashi et caractériser un cercle
\end{itemize}\end{frame}
\renewcommand{\slidekind}{COMPRENDRE}
\begin{frame}[t,fragile]{Une force et un déplacement}
Une force $\vec F$ de norme $10$ et un déplacement $\vec d$ de norme $3$ forment un angle de $60^\circ$. Le nombre $10\times3\times\cos60^\circ=15$ mesure la contribution de la force dans la direction du déplacement.\medskip\par Ce calcul associe aux deux vecteurs un \textbf{nombre réel}, appelé produit scalaire. Il tient compte à la fois des longueurs et de leur direction.
\end{frame}
\renewcommand{\slidekind}{COMPRENDRE}
\begin{frame}[t,fragile]{Définition par un angle}
Pour deux vecteurs non nuls,\[\vec u\cdot\vec v=\|\vec u\|\,\|\vec v\|\cos\theta,\qquad\theta\in[0,\pi].\]Si l’un des vecteurs est nul, le produit scalaire vaut $0$.\medskip\par Le résultat est un réel, jamais un vecteur. L’angle est celui de deux représentants de même origine.
\end{frame}
\renewcommand{\slidekind}{COMPRENDRE}
\begin{frame}[t,fragile]{Le signe donne une information}
Les normes sont positives. Le signe dépend donc de $\cos\theta$.\begin{itemize}\item Angle aigu : produit positif.\item Angle droit : produit nul.\item Angle obtus : produit négatif.\end{itemize}Deux vecteurs de même sens ont un produit positif. Deux vecteurs de sens opposés ont un produit négatif.
\end{frame}
\renewcommand{\slidekind}{COMPRENDRE}
\begin{frame}[t,fragile]{Projection orthogonale}
Soit $H$ le projeté orthogonal de $B$ sur la droite $(OA)$, avec $A\ne O$.\[\overrightarrow{OA}\cdot\overrightarrow{OB}=\overrightarrow{OA}\cdot\overrightarrow{OH}.\]Si $\overrightarrow{OH}=t\overrightarrow{OA}$, ce produit vaut $t\,OA^2$. Le réel $t$ peut être négatif : une projection située de l’autre côté de $O$ donne un produit négatif.
\end{frame}
\renewcommand{\slidekind}{MÉTHODE \& EXEMPLE}
\begin{frame}[t,fragile]{Exemple par projection}
On a $OA=5$. Le projeté $H$ de $B$ est sur la demi-droite $[OA)$ et $OH=2$. Alors\[\overrightarrow{OA}\cdot\overrightarrow{OB}=5\times2=10.\]Si $H$ était sur la demi-droite opposée, avec la même distance $OH=2$, le produit vaudrait $-10$. La longueur $OB$ n’est pas nécessaire ici.
\end{frame}
\renewcommand{\slidekind}{COMPRENDRE}
\begin{frame}[t,fragile]{Expression en coordonnées}
Dans un \textbf{repère orthonormé}, si $\vec u=(x,y)$ et $\vec v=(x\prime,y\prime)$, alors\[\boxed{\vec u\cdot\vec v=xx\prime+yy\prime}.\]Exemple : $(2,-3)\cdot(4,1)=2\times4-3\times1=5$.\medskip\par Cette formule ne s’utilise pas sans vérifier que les axes sont perpendiculaires et de même unité.
\end{frame}
\renewcommand{\slidekind}{COMPRENDRE}
\begin{frame}[t,fragile]{Symétrie et bilinéarité}
Pour tous vecteurs et tout réel $\lambda$ :\[\vec u\cdot\vec v=\vec v\cdot\vec u,\]\[\vec u\cdot(\vec v+\vec w)=\vec u\cdot\vec v+\vec u\cdot\vec w,\]\[\vec u\cdot(\lambda\vec v)=\lambda(\vec u\cdot\vec v).\]On développe donc comme une expression algébrique. Attention : un produit scalaire n’est pas un produit de trois vecteurs.
\end{frame}
\renewcommand{\slidekind}{JUSTIFIER}
\begin{frame}[t,fragile]{Pourquoi la formule fonctionne}
Dans une base orthonormée $(\vec i,\vec j)$, on écrit $\vec u=x\vec i+y\vec j$ et $\vec v=x\prime\vec i+y\prime\vec j$.\medskip\par La bilinéarité donne quatre termes. Comme $\vec i\cdot\vec i=\vec j\cdot\vec j=1$ et $\vec i\cdot\vec j=0$, les termes croisés disparaissent :\[\vec u\cdot\vec v=xx\prime+yy\prime.\]
\end{frame}
\renewcommand{\slidekind}{COMPRENDRE}
\begin{frame}[t,fragile]{Norme et carré scalaire}
\[\vec u\cdot\vec u=\|\vec u\|^2.\]Ainsi, pour $\vec u=(3,-4)$, $\|\vec u\|=\sqrt{9+16}=5$.\medskip\par En développant :\[\|\vec u+\vec v\|^2=\|\vec u\|^2+2\vec u\cdot\vec v+\|\vec v\|^2.\]Cette identité relie un calcul vectoriel aux longueurs.
\end{frame}
\renewcommand{\slidekind}{COMPRENDRE}
\begin{frame}[t,fragile]{Retrouver le produit avec trois longueurs}
Le développement de $\|\vec u-\vec v\|^2$ donne\[\vec u\cdot\vec v=\frac{\|\vec u\|^2+\|\vec v\|^2-\|\vec u-\vec v\|^2}{2}.\]Dans un triangle $ABC$, avec $\vec u=\overrightarrow{AB}$ et $\vec v=\overrightarrow{AC}$, on obtient\[\overrightarrow{AB}\cdot\overrightarrow{AC}=\frac{AB^2+AC^2-BC^2}{2}.\]
\end{frame}
\renewcommand{\slidekind}{MÉTHODE \& EXEMPLE}
\begin{frame}[t,fragile]{Orthogonalité : un critère de calcul}
Deux vecteurs non nuls sont orthogonaux si et seulement si leur produit scalaire est nul.\medskip\par Exemple : $(2,3)\cdot(3,-2)=6-6=0$.\medskip\par Pour prouver que $(AB)\perp(AC)$, vérifier que $A\ne B$, $A\ne C$, puis calculer $\overrightarrow{AB}\cdot\overrightarrow{AC}$. Le vecteur nul est orthogonal à tout vecteur, mais ne dirige aucune droite.
\end{frame}
\renewcommand{\slidekind}{COMPRENDRE}
\begin{frame}[t,fragile]{Chercher une coordonnée}
On veut $\vec u=(3,-2)$ orthogonal à $\vec v=(t,6)$.\[\vec u\cdot\vec v=3t-12.\]L’orthogonalité équivaut à $3t-12=0$, donc $t=4$.\medskip\par Vérification : $(3,-2)\cdot(4,6)=12-12=0$. L’équation obtenue est linéaire en l’inconnue.
\end{frame}
\renewcommand{\slidekind}{MÉTHODE \& EXEMPLE}
\begin{frame}[t,fragile]{Calculer un angle}
Pour $\vec u=(1,0)$ et $\vec v=(1,\sqrt3)$ :\[\vec u\cdot\vec v=1,\quad\|\vec u\|=1,\quad\|\vec v\|=2.\]Donc $\cos\theta=1/2$ et $\theta=\pi/3=60^\circ$.\medskip\par En général, calculer d’abord le produit et les deux normes, puis utiliser $\arccos$ dans l’unité demandée.
\end{frame}
\renewcommand{\slidekind}{COMPRENDRE}
\begin{frame}[t,fragile]{Théorème d’Al-Kashi}
Dans tout triangle $ABC$,\[\boxed{BC^2=AB^2+AC^2-2AB\times AC\cos\widehat{BAC}}.\]C’est une généralisation de Pythagore. Si l’angle en $A$ est droit, son cosinus est nul, et l’on retrouve $BC^2=AB^2+AC^2$. On utilise ce résultat avec deux côtés et l’angle compris entre eux.
\end{frame}
\renewcommand{\slidekind}{JUSTIFIER}
\begin{frame}[t,fragile]{Preuve d’Al-Kashi}
On a $\overrightarrow{BC}=\overrightarrow{AC}-\overrightarrow{AB}$. Par l’identité de norme,\[BC^2=AC^2+AB^2-2\overrightarrow{AB}\cdot\overrightarrow{AC}.\]Puis la définition par le cosinus donne\[\overrightarrow{AB}\cdot\overrightarrow{AC}=AB\times AC\cos\widehat{BAC}.\]La substitution donne exactement la formule annoncée.
\end{frame}
\renewcommand{\slidekind}{COMPRENDRE}
\begin{frame}[t,fragile]{Une longueur dans un triangle}
Si $AB=5$, $AC=8$ et $\widehat{BAC}=60^\circ$,\[BC^2=25+64-2\times5\times8\times\frac12=49.\]Comme une longueur est positive, $BC=7$.\medskip\par Le signe devant le terme contenant le cosinus est négatif. Si l’angle est obtus, son cosinus est négatif et ce terme augmente le carré du côté opposé.
\end{frame}
\renewcommand{\slidekind}{COMPRENDRE}
\begin{frame}[t,fragile]{Cercle de diamètre donné}
Pour $A\ne B$, le cercle de diamètre $[AB]$ est l’ensemble des points $M$ tels que\[\overrightarrow{MA}\cdot\overrightarrow{MB}=0.\]Pour $M\ne A,B$, cela signifie que le triangle $AMB$ est rectangle en $M$. Les points $A$ et $B$ satisfont aussi l’égalité grâce au vecteur nul.
\end{frame}
\renewcommand{\slidekind}{JUSTIFIER}
\begin{frame}[t,fragile]{Preuve de la caractérisation du cercle}
Soit $I$ le milieu de $[AB]$. On a $\overrightarrow{IA}=-\overrightarrow{IB}$. Alors\[\overrightarrow{MA}\cdot\overrightarrow{MB}=(\overrightarrow{MI}+\overrightarrow{IA})\cdot(\overrightarrow{MI}-\overrightarrow{IA})=MI^2-IA^2.\]Le produit est nul exactement lorsque $MI=IA$, c’est-à-dire lorsque $M$ est sur le cercle de centre $I$ et de rayon $AB/2$.
\end{frame}
\renewcommand{\slidekind}{MÉTHODE \& EXEMPLE}
\begin{frame}[t,fragile]{Quelle formule choisir ?}
\begin{itemize}\item Coordonnées connues : $xx\prime+yy\prime$.\item Normes et angle connus : $\|u\|\|v\|\cos\theta$.\item Trois longueurs connues : identité de norme.\item Projection simple : utiliser la longueur algébrique projetée.\end{itemize}Toujours dessiner les vecteurs de même origine avant de lire l’angle. $\overrightarrow{BA}=-\overrightarrow{AB}$ change le signe du produit.
\end{frame}
\renewcommand{\slidekind}{À RETENIR}
\begin{frame}[t,fragile]{Bilan et contrôle rapide}
Pour $AB=AC=5$ et $BC=6$,\[\overrightarrow{AB}\cdot\overrightarrow{AC}=\frac{25+25-36}{2}=7.\]L’angle en $A$ est aigu car le produit est positif. Son cosinus vaut $7/25$.\medskip\par À retenir : un produit scalaire transforme une question géométrique en un calcul de réel. La précision des orientations et du repère est essentielle.
\end{frame}
\renewcommand{\slidekind}{COMPRENDRE}
\begin{frame}[t,fragile]{Des coordonnées à l’angle}
$\vec u=(1,2)$ et $\vec v=(3,-1)$.\[\vec u\cdot\vec v=1\times3+2\times(-1)=1.\]\[\|\vec u\|=\sqrt5,\quad\|\vec v\|=\sqrt{10},\quad\cos\theta=\frac1{\sqrt{50}}.\] Le produit est positif : l’angle non orienté est aigu. La calculatrice donne environ $81,9^\circ$.
\end{frame}
\renewcommand{\slidekind}{MÉTHODE \& EXEMPLE}
\begin{frame}[t,fragile]{Un calcul par les longueurs}
Un triangle a $AB=4$, $AC=5$ et $BC=3$.\[\overrightarrow{AB}\cdot\overrightarrow{AC}=\frac{AB^2+AC^2-BC^2}{2}=\frac{16+25-9}{2}=16.\] Donc $\cos\widehat{BAC}=16/(4\times5)=0,8$. Les deux vecteurs doivent partir du même sommet pour appliquer cette lecture de l’angle.
\end{frame}
\renewcommand{\slidekind}{COMPRENDRE}
\begin{frame}[t,fragile]{Attention à l’orientation des vecteurs}
$\overrightarrow{BA}=-\overrightarrow{AB}$. Par bilinéarité,\[\overrightarrow{BA}\cdot\overrightarrow{AC}=-\overrightarrow{AB}\cdot\overrightarrow{AC}.\] Inverser un seul vecteur change le signe du produit. Inverser les deux conserve le produit. Un dessin aide à vérifier les directions avant de remplacer par des longueurs.
\end{frame}
\renewcommand{\slidekind}{JUSTIFIER}
\begin{frame}[t,fragile]{Un critère d’orthogonalité justifié}
$\vec u=(2,-3)$ et $\vec v=(t,4)$. Ils sont orthogonaux si et seulement si\[2t-12=0\quad\Longleftrightarrow\quad t=6.\] Pour $t=6$, le produit vaut bien $0$. La conclusion porte sur les directions des vecteurs, sans exiger qu’ils soient tracés avec la même origine.
\end{frame}
\renewcommand{\slidekind}{COMPRENDRE}
\begin{frame}[t,fragile]{Le produit scalaire par projection}
\begin{columns}[c,onlytextwidth]\begin{column}{.60\textwidth}\centering\begin{adjustbox}{max width=\linewidth,max height=3.7cm}\begin{tikzpicture}[>=Stealth]\draw[->,thick,Blue](0,0)--(4,0)node[below]{$\vec u$};\draw[->,thick,Ink](0,0)--(2.5,2)node[above]{$\vec v$};\draw[dashed](2.5,2)--(2.5,0)node[below]{$H$};\draw(2.25,0)--(2.25,.25)--(2.5,.25);\node[below left]at(0,0){$O$};\draw[Blue](.7,0)arc(0:38.66:.7);\node at(1,.3){$\theta$};\end{tikzpicture}\end{adjustbox}\end{column}\begin{column}{.36\textwidth} La projection de $\vec v$ sur la direction de $\vec u$ a pour longueur $OH=\|\vec v\|\cos\theta$. Ici l’angle est aigu, donc le produit est positif.\end{column}\end{columns}
\end{frame}
\renewcommand{\slidekind}{À VOUS DE JOUER}\begin{frame}[t]{Vérification des acquis 1}
Si $\|u\|=3$, $\|v\|=4$ et leur angle vaut $60^\circ$, alors $u\cdot v=$\begin{enumerate}
\item $12$
\item $6$
\item $0$
\item $-6$
\end{enumerate}\medskip Une seule réponse est correcte. Justifier votre choix.\end{frame}
\renewcommand{\slidekind}{COMPRENDRE LA RÉPONSE}\begin{frame}[t]{Vérification 1 : résolution}
\begin{block}{Réponse 2}$6$\end{block}\medskip $3\times4\times\cos60^\circ=12\times\frac12=6$.\end{frame}
\renewcommand{\slidekind}{À VOUS DE JOUER}\begin{frame}[t]{Vérification des acquis 2}
Pour $A\ne B$, l’ensemble $\overrightarrow{MA}\cdot\overrightarrow{MB}=0$ est…\begin{enumerate}
\item la médiatrice de $[AB]$
\item la droite $(AB)$
\item le disque de diamètre $[AB]$
\item le cercle de diamètre $[AB]$
\end{enumerate}\medskip Une seule réponse est correcte. Justifier votre choix.\end{frame}
\renewcommand{\slidekind}{COMPRENDRE LA RÉPONSE}\begin{frame}[t]{Vérification 2 : résolution}
\begin{block}{Réponse 4}le cercle de diamètre $[AB]$\end{block}\medskip En notant $I$ le milieu, le produit vaut $MI^2-IA^2$. Il est nul exactement sur le cercle.\end{frame}
\end{document}
